\section{Oracle-Bounded Recoverable Envelope (R2-04)}
\label{sec:recoverable-envelope}

R2-04 synthesizes the R2-02/R2-03 evidence into a single decomposition,
applied at every tested budget, and reported here as the headline II-A
result. No new replay, model, or runtime invocation is performed in this
step; it is a pure offline synthesis over already-registered evidence.

\begin{definition}[Structural / information / total decomposition]
\label{def:decomposition}
Let $S$, $G$, $O$ denote the miss-byte totals (at a fixed budget
$\budget$) of Runtime-Sharded-LRU, Global-LRU, and Oracle-Global-Belady
respectively. Define
\begin{align}
\text{Structural sharding gap} &= S - G, \label{eq:structural} \\
\text{Policy-information headroom} &= G - O, \label{eq:information} \\
\text{Total observed upper-bound gap} &= S - O. \label{eq:total}
\end{align}
By construction, Eq.~\eqref{eq:structural} $+$ Eq.~\eqref{eq:information}
$=$ Eq.~\eqref{eq:total} exactly, for both the miss-byte and miss-count
views, independent of the sign of the structural term. The corresponding
\emph{relative} fractions (each normalized by its own comparison's first
operand) are \textbf{not} summed arithmetically, since their denominators
differ -- see Table~\ref{tab:r2-03-sweep} vs.\ the structural
percentages quoted in Section~\ref{sec:sensitivity-capacity}.
\end{definition}

The \emph{policy-information headroom} (Eq.~\ref{eq:information}) is the
quantity of scientific interest for motivating II-B/II-C: it isolates the
cost of imperfect \emph{causal information}, independent of the specific
sharded retention architecture. The \emph{structural sharding gap}
(Eq.~\ref{eq:structural}) is reported for completeness but is not the
focus of this report, and carries the limitation that the offline
Runtime-Sharded-LRU's shard-to-key mapping is statistical rather than the
runtime's exact address-derived mapping.

\begin{table}[t]
\centering
\caption{Headline II-A synthesis (R2-04): oracle-bounded recoverable-envelope summary.}
\label{tab:r2-04-synthesis}
\small
\begin{tabular}{@{}lr@{}}
\toprule
Quantity & Value \\
\midrule
Information headroom at 4~GiB (anchor) & 58.59\% \\
Maximum tested relative information headroom & 70.66\% (at 8~GiB) \\
Maximum tested absolute miss-byte headroom (GB) & 103.57 (at 1.5~GiB) \\
First tested policy-indifference budget (GiB) & 16 \\
Oracle strictly improves at every capacity-constrained point & yes \\
Observed information-headroom curve shape & collapse at effective residency \\
Causally recovered margin & \textsc{not yet measured} (\S~II-C) \\
\bottomrule
\end{tabular}
\end{table}


\paragraph{Absolute vs.\ relative headroom occur at different budgets.}
A single important distinction this synthesis makes explicit: the budget
at which the \emph{relative} information headroom is maximized ($8$~GiB,
$70.66\%$) is \textbf{not} the budget at which the \emph{absolute}
miss-byte headroom is maximized ($1.5$~GiB, $103{,}574{,}900{,}736$
bytes $\approx 103.57$~GB). Under more severe memory pressure (lower
budgets), both Global-LRU and the Oracle incur far more total misses, so
the \emph{absolute} volume of miss bytes the Oracle avoids can be larger
even while its \emph{relative} advantage (as a fraction of Global-LRU's
own, larger, miss-byte total) is smaller. Neither framing is more
``correct''; they answer different engineering questions (fraction of
traffic avoidable vs.\ total bytes avoidable) and are reported
separately, never conflated into one number.

\paragraph{Headroom persistence under pressure.} At every one of the
four capacity-constrained tested budgets ($1.5$--$8$~GiB, where
Global-LRU's own eviction count is nonzero), the Oracle strictly
dominates Global-LRU: a strictly positive miss and miss-byte reduction
is observed at each such point. The $16$~GiB
policy-indifference point is excluded from this statement by definition
(no policy has anything left to decide there, so ``the oracle strictly
improves'' is not a meaningful comparison at that point).

\paragraph{Observed curve shape.} The relative information-headroom
curve (Figure~\ref{fig:headroom-vs-budget}) is classified, deterministically
from the five tested points, as a \emph{rise followed by a collapse at
the observed effective-residency budget}: non-decreasing across
$1.5$--$8$~GiB, then dropping to its global minimum ($0$) exactly at
$16$~GiB, the first tested budget exhibiting policy convergence. This
classification is derived mechanically from the tested points, not
assumed in advance.

\begin{definition}[Future II-C metric -- documented, not evaluated]
\label{def:ii-c-metric}
For a causal baseline with miss bytes $B_\budget$ and a future causal
predictor-driven policy with miss bytes $P_\budget$, both at budget
$\budget$, define
\begin{equation}
\mathrm{FractionOfOracleHeadroomRecovered}_\budget \;=\;
\frac{B_\budget - P_\budget}{B_\budget - O_\budget}, \label{eq:ii-c}
\end{equation}
where $O_\budget$ is the Oracle's miss bytes at that budget. This metric
is defined now, as the scientific bridge from II-A's upper bound to
II-C's demonstrated causal recovery, but \textbf{cannot be evaluated} --
no causal predictor $P$ exists yet. It is reported here purely as a
predeclared target metric for Research II-C.
\end{definition}

\paragraph{Causal recovery: explicitly unmeasured.} We emphasize, as
strongly as the evidence allows: nothing in this section measures or
estimates the fraction of the policy-information headroom that any
practical causal predictor could actually recover. That quantity is
\notyetmeasured{} and is the subject of Research II-C
(Section~\ref{sec:iic-placeholder}), which has not started.
